% \documentclass[a4paper,12pt]{article}{{{
\documentclass[a4paper,11pt]{article}
% -------------------------------------------------
% Set up the page margins.
% -------------------------------------------------
% \usepackage[text={6.5in,9.5in},centering]{geometry}
% \setlength{\topmargin}{-0.25in}
\usepackage[margin=1in]{geometry}

% -------------------------------------------------\left( \left[ \frac{a-d}{2}\right]q + p \right)
% Load Packages here
% -------------------------------------------------
%\usepackage{hyperref}
\usepackage{graphicx,url,subfig}
\RequirePackage[OT1]{fontenc}
\RequirePackage{amsthm,amsmath,amssymb,amscd}
\RequirePackage[numbers]{natbib}
\RequirePackage[colorlinks,citecolor=blue,urlcolor=blue]{hyperref}
\usepackage{graphics}
\usepackage{enumerate}
\usepackage{datetime}
\usepackage{etex}
% \usepackage[notcite,notref,color]{showkeys}
\usepackage[normalem]{ulem} % either use this (simple) or
\usepackage{soul} % use this (many fancier options)
\makeatother
\numberwithin{equation}{section}
\allowdisplaybreaks
\usepackage{mathrsfs}


% -------------------------------------------------
% check references
% -------------------------------------------------
% \usepackage{refcheck}
% \usepackage[notref,notcite]{showkeys}
% \usepackage[notcite]{showkeys}
% \usepackage{showkeys}

% -------------------------------------------------
% Environments here
% -------------------------------------------------
\theoremstyle{plain}
\newtheorem{theorem}{Theorem}[section]
\newtheorem{corollary}[theorem]{Corollary}
\newtheorem{lemma}[theorem]{Lemma}
\newtheorem{proposition}[theorem]{Proposition}
\newtheorem{exercise}{Exercise}
\theoremstyle{definition}
\newtheorem{definition}[theorem]{Definition}
\newtheorem{hypothesis}[theorem]{Hypothesis}
\newtheorem{assumption}[theorem]{Assumption}
\newtheorem{question}{Question}

\newtheorem{remark}[theorem]{Remark}
\newtheorem{conjecture}[theorem]{Conjecture}
\newtheorem{example}[theorem]{Example}

% -------------------------------------------------
%  Define short hand symbols.
% -------------------------------------------------
\newcommand{\E}{\mathbb{E}}
\newcommand{\W}{\dot{W}}
\newcommand{\B}{\textbf{B}}
\newcommand{\D}{\mathbb{D}}
\newcommand{\ud}{\ensuremath{\mathrm{d} }}
\newcommand{\Ceil}[1]{\left\lceil #1 \right\rceil}
\newcommand{\Floor}[1]{\left\lfloor #1 \right\rfloor}
\newcommand{\sgn}{\text{sgn}}
\newcommand{\Lad}{\text{L}_{\text{ad}}^2}
\newcommand{\SI}[1]{\mathcal{I}\left[#1 \right]}
\newcommand{\SIB}[2]{\mathcal{I}_{#2}\left[#1 \right]}
\newcommand{\Indt}[1]{1_{\left\{#1 \right\} }}
\newcommand{\LadInPrd}[1]{\left\langle #1 \right\rangle_{\text{L}_\text{ad}^2}}
\newcommand{\LadNorm}[1]{\left|\left|  #1 \right|\right|_{\text{L}_\text{ad}^2}}
\newcommand{\Norm}[1]{\left\|  #1   \right\|}
\newcommand{\Ito}{It\^{o} }
\newcommand{\Itos}{It\^{o}'s }
\newcommand{\spt}[1]{\text{supp}\left(#1\right)}
\newcommand{\InPrd}[1]{\left\langle #1 \right\rangle}
\newcommand{\mr}{\textbf{r}}
\newcommand{\Ei}{\text{Ei}}
\newcommand{\arctanh}{\operatorname{arctanh}}
\newcommand{\ind}[1]{\mathbb{I}_{\left\{ {#1} \right\} }}

\DeclareMathOperator{\esssup}{\ensuremath{ess\,sup}}

\newcommand{\steps}[1]{\vskip 0.3cm \textbf{#1}}
\newcommand{\calB}{\mathcal{B}}
\newcommand{\calD}{\mathcal{D}}
\newcommand{\calE}{\mathcal{E}}
\newcommand{\calF}{\mathcal{F}}
\newcommand{\calG}{\mathcal{G}}
\newcommand{\calK}{\mathcal{K}}
\newcommand{\calH}{\mathcal{H}}
\newcommand{\calI}{\mathcal{I}}
\newcommand{\calL}{\mathcal{L}}
\newcommand{\calM}{\mathcal{M}}
\newcommand{\calN}{\mathcal{N}}
\newcommand{\calO}{\mathcal{O}}
\newcommand{\calT}{\mathcal{T}}
\newcommand{\calP}{\mathcal{P}}
\newcommand{\calR}{\mathcal{R}}
\newcommand{\calS}{\mathcal{S}}
\newcommand{\bbC}{\mathbb{C}}
\newcommand{\bbN}{\mathbb{N}}
\newcommand{\bbP}{\mathbb{P}}
\newcommand{\bbZ}{\mathbb{Z}}
\newcommand{\myVec}[1]{\overrightarrow{#1}}
\newcommand{\sincos}{\begin{array}{c} \cos \\ \sin \end{array}\!\!}
\newcommand{\CvBc}[1]{\left\{\:#1\:\right\}}
\newcommand*{\one}{ {{\rm 1\mkern-1.5mu}\!{\rm I} }}
\newcommand{\RR}{\mathbb{R}}
\newcommand{\R}{\mathbb{R}}
\newcommand{\myEnd}{\hfill$\square$}
\newcommand{\ds}{\displaystyle}
\newcommand{\Shi}{\text{Shi}}
\newcommand{\Chi}{\text{Chi}}
\newcommand{\Daw}{\ensuremath{\mathrm{daw} }}
\newcommand{\Erf}{\ensuremath{\mathrm{erf} }}
\newcommand{\Erfi}{\ensuremath{\mathrm{erfi} }}
\newcommand{\Erfc}{\ensuremath{\mathrm{erfc} }}
\newcommand{\He}{\ensuremath{\mathrm{He} }}

\def\crtL{\theta}

\newcommand{\conRef}[1]{(#1)}

\DeclareMathOperator{\Lip}{\mathit{L}}
\DeclareMathOperator{\LIP}{Lip}
\DeclareMathOperator{\lip}{\mathit{l}}

\DeclareMathOperator{\Vip}{\overline{\varsigma}}
\DeclareMathOperator{\vip}{\underline{\varsigma}}
\DeclareMathOperator{\vv}{\varsigma}

\usepackage[notref,notcite]{showkeys}

\title{Nick Eisenberg's Meeting Notes}
\author{Le Chen, Nick Eisenberg, Zhijian Wu}
\date{October 2019}

\usepackage{natbib}
\usepackage{graphicx}
% }}}
\begin{document}

\begin{question}[Norm increments in time and space.]
	If we want to keep the more general initial conditions $u_0$ and $v_0$ then we need to find an upper bound of the following form:
	\[
		\int_{\R}|G(s,z-y) - G(t,x-y)| \ud y \le C(T)(|s-t|+|z-x|) \quad \text{for } t \in [0,T] \text{ and } x,z \in \R.
	\]
	To simplify this at first we can look at both 
	\begin{align}
		\label{E:HGspace} & \int_{\R}|G(t,z-y) - G(t,x-y)| \ud y \\
		\label{E:HGtime} & \int_{\R}|G(s, x-y) - G(t,x-y)| \ud y 
	\end{align}
	and in addition consider the Heat equation case,
	\[
		G(t,x) = \frac{1}{\sqrt{2 \pi t}}\exp\left( - \frac{x^2}{2t}\right).
	\]
	Starting with \eqref{E:HGspace}, we first assume that $x < z$ and use the Mean Value Theorem to get 
	\begin{align*}
		\int_{\R}|G(t,z-y) - G(t,x-y)| \ud y & \le  |z - x| 	\int_{\R}\left|\frac{\partial G}{\partial x}(t,\xi(x,z) -y)\right|\ud y 
	\end{align*}
	where $\xi(x, z) \in (x,z)$ is a constant. Using Mathematica,
	\begin{verbatim}
		g[x_] := (2 Pi t)^(-1/2) Exp[-(x^2)/(2t)] 
		
		Integrate[Abs[D[g[x-y],x]],{y,-Infinity,Infinity}, Assumptions -> t>0 
		&& -Infinity < x < Infinity]
	\end{verbatim}
	we get that for any $x\in \R$ that 
	\[
		\int_{\R}\left|\frac{\partial G}{\partial x}(t,x-y)\right|\ud y = \frac{\sqrt{\frac{2}{\pi}}}{\sqrt{t}}
	\]
	and hence 
	\begin{align*}
		\int_{\R}|G(t,z-y) - G(t,x-y)| \ud y & \le \sqrt{\frac{2}{\pi}} \frac{1}{\sqrt{t}}|z - x|.
	\end{align*}
	Similarly for \eqref{E:HGtime}, by assuming $t<s$ we get 
	\begin{align*}
	\int_{\R}|G(s,x-y) - G(t,x-y)| \ud y & \le  |s - t| 	\int_{\R}\left|\frac{\partial G}{\partial t}(\eta(t,s), x -y)\right|\ud y 
	\end{align*}
	where $\eta(t,s) \in (t, s)$ is a constant. And by using Mathematica, 
	\begin{verbatim}
		Integrate [Abs[D[g[x-y],t]],{y,-Infinity, Infinity}, Assumptions -> t>0 
		&& -Infinity < x < Infinity]
	\end{verbatim}
	we get that for $t>0$ and $x \in \R$ that 
	\begin{align*}
		\int_{\R}\left|\frac{\partial G}{\partial t}(t, x -y)\right|\ud y  & = \sqrt{\frac{2}{e \pi}} \frac{1}{t}
	\end{align*}
	and therefore 
	\begin{align*}
	\int_{\R}|G(s,x-y) - G(t,x-y)| \ud y & \le \sqrt{\frac{2}{e \pi}} \frac{1}{\eta(t,s)}  |s - t| .
	\end{align*}
	Now we can see that 
	\begin{align*}
	\int_{\R}|G(s,z-y) - G(t,x-y)| \ud y &\le  \int_{\R}|G(t,z-y) - G(t,x-y)| \ud y +  \int_{\R}|G(s,z-y) - G(t,z-y)| \ud y  \\
	& \le \sqrt{\frac{2}{\pi}} \frac{1}{\sqrt{t}}|z - x| + \sqrt{\frac{2}{e \pi}} \frac{1}{\eta(t,s)}  |s - t|.
	\end{align*}
	The $t$ and $\eta(t,s)$ seem to be an issue since they are unbounded as $t, s \to 0$. 
	
	\begin{remark}
		We can tweak the upper bound for \eqref{E:HGtime} at little to remove the $\eta(t,s)$. First note that 
			\begin{align*}
			\left|\frac{\partial G}{\partial t}(t, x -y)\right| & = \left| -\frac{\exp\left(-\frac{x^2}{2t}\right)}{2\sqrt{2 \pi} t^{3/2}} + \frac{\exp\left(-\frac{x^2}{2t}\right) x^2}{2\sqrt{2 \pi} t^{5/2} }\right| \\
			& \le \frac{\exp\left(-\frac{x^2}{2t}\right)}{2\sqrt{2 \pi} t^{3/2}}  + \frac{\exp\left(-\frac{x^2}{2t}\right) x^2}{2\sqrt{2 \pi} t^{5/2} }
			\end{align*}
		and notice that for $t < s$ we have that 
		\begin{align*}
			\frac{\exp\left(-\frac{x^2}{2t}\right)}{2\sqrt{2 \pi} t^{3/2}}  + \frac{\exp\left(-\frac{x^2}{2t}\right) x^2}{2\sqrt{2 \pi} t^{5/2} } & \le
			 \frac{\exp\left(-\frac{x^2}{2s}\right)}{2\sqrt{2 \pi} t^{3/2}}  + \frac{\exp\left(-\frac{x^2}{2s}\right) x^2}{2\sqrt{2 \pi} t^{5/2} }
		\end{align*}
		and therefore 
		\begin{align*}
			\int_{\R}\left|\frac{\partial G}{\partial t}(t, x -y)\right|\ud y  & \le 		 \int_{\R} \frac{\exp\left(-\frac{(x-y)^2}{2s}\right)}{2\sqrt{2 \pi} t^{3/2}}  \ud y+ \int_{\R} \frac{\exp\left(-\frac{(x-y)^2}{2s}\right) x^2}{2\sqrt{2 \pi} t^{5/2} } \ud y \\
			&\le \frac{1}{2t}\sqrt{\frac{s}{t}} + \frac{s}{2t^2}\sqrt{\frac{s}{t}}
		\end{align*}
		which implies that 
		\[
				\int_{\R}|G(s,x-y) - G(t,x-y)| \ud y \le \left(\frac{1}{2t}\sqrt{\frac{s}{t}} + \frac{s}{2t^2}\sqrt{\frac{s}{t}}\right)|s-t|.
		\]
	However we still run into the same issue. 
	\end{remark}
		
\end{question}
	
	
	
	\addcontentsline{toc}{section}{Bibliography}
	\begin{thebibliography}{999}
		%sample
		%\bibitem{Chen17Tran}% {{{ 3.
		%Chen, L. (2017).
		%\newblock Nonlinear stochastic time-fractional diffusion equations on $\R$: moments, H\"older regularity and intermittency.
		%\newblock {\em Trans. Amer. Math. Soc.} \textbf{369} (12) 8497--8535.
		
		\bibitem{CHN19}% {{{ 5.
		Chen, L., Hu, Y., and Nualart, D. (2019).
		\newblock Nonlinear stochastic time-fractional slow and fast diffusion equations on $\R^d$.
		\newblock {\em Stochastic Process. Appl.} \textbf{129}, 5073--5112
		% }}}
		
		\bibitem{SoleMillet2021}
		Sanz-Sole, M. and Millet, A. (2021)
		\newblock{\em Global solutions to stochastic wave equations with superlinear coefficients.}
		\newblock Stochastic Processes and their Applications, Elsevier, 2021, 139, pp.175-211.
	\end{thebibliography}
	
\end{document}
